% ============================================================
% Second-Shell-to-Second-Shell Edge-Direction Perpendicularity (G2.4)
%   for the 600-Cell Edge Graph at Vertex-Aligned Reading C
% Publication-Grade Hardening of the Second-Shell Analog of
%   the First-Shell Edge-Direction Perpendicularity (Patch 0551)
% Conscious Point Physics Flagship Paper Series
%   (F.1 Dynamical Substrate Law sub-question hardened-theorem artifact;
%    EIGHTH artifact in F.1 Layer 3 hardened-theorem sequence;
%    THIRD artifact in OPEN-FP-F1-1 Route C sequence 2 after Patch 0587
%    G2 hardening; parallel partner to Patch 0588 host-to-second-shell
%    projection G2.1. Hardens G2.4 -- the perpendicularity of all 30
%    second-shell-to-second-shell 600-cell edges to n_hat at vertex-
%    aligned Reading C -- at publication-grade Layer 3 UNCONDITIONAL
%    rigor via G2 hardened at Patch 0587 + standard chord-difference
%    argument.)
%
% Patch 0589 (Session 145, 26 May 2026) -- third OPEN-FP-F1-1 Route
%   C sequence 2 closure Patch; parallel partner to Patch 0588
%   (host-to-second-shell projection G2.1). Patch number skips 0586
%   (reserved for concurrent Opus window's Session 144 handover
%   completion documentation).
% ============================================================
% Version 1.2 --- 17 August 2026 (Patch 3214):
%   added a How-we-review methodology note stating plainly that AI panel
%   review is a self-applied verification method and not journal peer
%   review; twelve chirality papers retitled so the descriptive title
%   leads and the identifier is demoted to a subtitle. No claim,
%   equation, number or section changed.
% Version 1.1 --- 17 August 2026 (Patch 3213):
%   extended the generated identifier appendix to all five identifier
%   families (THEO-, FI-, PRED-, CONJ-, PH- as well as OPEN-), and
%   replaced review-convergence scores with AI review panel wording
%   where present. No claim, equation, number or section changed.
% Version 1.0 --- 17 August 2026 (Patch 3212):
%   added the generated plain-language appendix glossing the internal
%   programme identifiers used in this paper (founder jargon ruling, 16
%   Aug 2026). No claim, equation, number or section changed.
%   This is the FIRST version stamp assigned to this paper; 1.0 denotes
%   the first stamped revision, NOT a claim of v1.0 shipped grade.

\documentclass[11pt,letterpaper]{article}
\usepackage[utf8]{inputenc}
\usepackage[margin=1in]{geometry}
\usepackage[T1]{fontenc}
\usepackage{amsmath,amssymb,amsfonts,amsthm}
\usepackage{xcolor}
\usepackage{hyperref}
\usepackage{enumitem}
\usepackage{parskip}
\usepackage{mdframed}

\hypersetup{
    colorlinks=true,
    linkcolor=blue,
    citecolor=blue,
    urlcolor=blue
}

\newtheorem{theorem}{Theorem}[section]
\newtheorem{proposition}[theorem]{Proposition}
\newtheorem{lemma}[theorem]{Lemma}
\newtheorem{corollary}[theorem]{Corollary}
\newtheorem{definition}[theorem]{Definition}
\newtheorem{remark}[theorem]{Remark}

\newcommand{\nhat}{\hat{n}}
\newcommand{\vhost}{v_{\text{host}}}
\newcommand{\Gsixcell}{\Gamma_{600}}
\newcommand{\Reals}{\mathbb{R}}
\newcommand{\Hthree}{H_3}
\newcommand{\Hfour}{H_4}
\newcommand{\Icoh}{I_h}

\title{Second-Shell-to-Second-Shell Edge-Direction Perpendicularity (G2.4) \\
in the 600-Cell at Vertex-Aligned Reading C \\
\large Publication-Grade Hardening of the Second-Shell Analog \\
of the Patch~0551 First-Shell Perpendicularity Theorem}

\author{Thomas Lee Abshier, ND, and Claude Opus 4.7 \\ Hyperphysics Institute --- F.1 Dynamical Substrate Law trajectory}

\date{26 May 2026 (Session 145, CPP Patch 0589)\\[2pt]
{\small Version~1.0 --- 17 August 2026 (Patch 3212: added the generated plain-language appendix glossing the internal programme identifiers used in this paper (founder jargon ruling, 16 Aug 2026). No claim, equation, number or section changed.)}\\[2pt]
{\small Version~1.1 --- 17 August 2026 (Patch 3213: extended the generated identifier appendix to all five identifier families (THEO-, FI-, PRED-, CONJ-, PH- as well as OPEN-), and replaced review-convergence scores with AI review panel wording where present. No claim, equation, number or section changed.)}\\[2pt]
{\small Version~1.2 --- 17 August 2026 (Patch 3214: added a How-we-review methodology note stating plainly that AI panel review is a self-applied verification method and not journal peer review; twelve chirality papers retitled so the descriptive title leads and the identifier is demoted to a subtitle. No claim, equation, number or section changed.)}}

\begin{document}
\maketitle

\begin{abstract}
We give a publication-grade Layer~3 \emph{unconditional} hardening of the second-shell-to-second-shell edge-direction perpendicularity identity G2.4 for the 600-cell at vertex-aligned Reading~C: the second-shell analog of the Patch~0551 first-shell-to-first-shell perpendicularity theorem. The hardened statement asserts that, under (H1)~vertex-aligned Reading~C with $\nhat = \vhost$, (H2)~600-cell unit-vertex normalisation, (H3)~the icosahedral residual symmetry $\Hthree = \Icoh$ at the host vertex, each of the 30 unit edge-direction vectors $\hat{e}_{j_1 j_2}$ of the second-shell-to-second-shell 600-cell edges (the 30 edges of the dodecahedron formed by $S_2(\vhost)$ under $\Icoh$) satisfies $\hat{e}_{j_1 j_2} \cdot \nhat = 0$ identically. The proof is essentially elementary: both endpoints of any second-shell-to-second-shell edge have $\nhat$-component equal to $1/2$ (by G2 hardened at Patch~0587), so their difference $w_{j_2} - w_{j_1}$ has $\nhat$-component zero, hence the edge-direction vector lies in the 3-plane $\nhat^\perp \subset \Reals^4$. The result requires no $I_h$-transitivity argument beyond what is implicit in the G2-uniformity inheritance (the perpendicularity holds for every pair $(w_{j_1}, w_{j_2})$ regardless of $\Icoh$-orbit structure). With Patch~0589 in hand, the F.1 Layer~3 hardened-theorem sequence advances to \textbf{eight artifacts}, completing the parallel pair of A3 + A5 Route~C sequence~2 artifacts unlocked by G2 hardening at Patch~0587. The remaining sequence-2 artifacts are A4 (first-shell-to-second-shell edge orbit classification G2.3, conditional on A2 + A3), A6 (path-class weight enumeration), and A7 (umbrella assembly producing the THEO-DSL-5 candidate $\alpha_2$).
\end{abstract}

\section{Introduction and statement of the main result}\label{sec:introduction}

This paper hardens the second-shell-to-second-shell edge-direction perpendicularity identity G2.4 at publication-grade Layer~3 unconditional rigor, executing the third artifact of Route~C sequence~2 of the OPEN-FP-F1-1 trajectory. The substantive content asserts that the 30 edges of the dodecahedron formed by $S_2(\vhost)$ (which coincide with the 30 second-shell-to-second-shell 600-cell edges; see Lemma~\ref{lem:30_edges_count}) have edge-direction vectors all perpendicular to $\nhat$ in the ambient $\Reals^4$.

\paragraph{Position in the F.1 Layer 3 hardening sequence.}
The present Patch~0589 is the \textbf{eighth} artifact in the F.1 Layer~3 hardened-theorem sequence and the \textbf{third} under Route~C sequence~2 of the OPEN-FP-F1-1 trajectory. The parallel partner artifact is Patch~0588 (host-to-second-shell projection G2.1), which is independent of the present perpendicularity result and was produced concurrently. Both Patches~0588 + 0589 were unlocked at publication-grade Layer~3 unconditional rigor by the Patch~0587 G2 hardening per Patch~0587 Corollary~4.1.

\paragraph{Statement of the main result.}
Let $\Gsixcell$ denote the 600-cell edge graph (Patch~0550 \S2.1; \cite{coxeter1973}). Fix a host vertex $\vhost \in V(\Gsixcell)$, and let $S_2(\vhost) = \{w_1, \ldots, w_{20}\}$ denote the second graph-distance shell of $\vhost$ in $\Gsixcell$ as characterised in Patch~0587 \S2.1. Let $E_{S_2}(\vhost) \subset E(\Gsixcell)$ denote the set of 600-cell edges with both endpoints in $S_2(\vhost)$.

\begin{lemma}[Edge count of $E_{S_2}(\vhost)$]\label{lem:30_edges_count}
$|E_{S_2}(\vhost)| = 30$, with each $w_j \in S_2(\vhost)$ having exactly 3 of its 12 600-cell-neighbours in $S_2(\vhost)$. The 30 edges of $E_{S_2}(\vhost)$ coincide with the 30 edges of the regular dodecahedron formed by the 20 second-shell vertices under the $\Icoh$ residual symmetry.
\end{lemma}

The proof of Lemma~\ref{lem:30_edges_count} is given in \S\ref{sec:edge_count_proof}. With this characterisation in hand, the main theorem is:

\begin{theorem}[Second-shell-to-second-shell edge-direction perpendicularity, hardened]\label{thm:g24_hardened}
Under hypotheses (H1)--(H3) of Patch~0587 Theorem~1.1, for every second-shell-to-second-shell 600-cell edge $(w_{j_1}, w_{j_2}) \in E_{S_2}(\vhost)$ with edge-direction unit vector
\begin{equation}\label{eq:edge_unit_def}
\hat{e}_{j_1 j_2} \;:=\; \frac{w_{j_2} - w_{j_1}}{|w_{j_2} - w_{j_1}|} \;\in\; \Reals^4,
\end{equation}
the projection on $\nhat$ vanishes identically:
\begin{equation}\label{eq:g24_main}
\hat{e}_{j_1 j_2} \cdot \nhat \;=\; 0.
\end{equation}
Equivalently, every second-shell-to-second-shell edge direction lies in the 3-plane $\nhat^\perp = \{x \in \Reals^4 : x \cdot \nhat = 0\}$ orthogonal to $\nhat$ in the ambient $\Reals^4$.
\end{theorem}

The proof (\S\ref{sec:proof}) is essentially a chord-difference calculation: both endpoints have $w_j \cdot \nhat = 1/2$ by G2 (Patch~0587 Theorem~1.1), so the difference $w_{j_2} - w_{j_1}$ has $\nhat$-component zero, hence the edge-direction has $\nhat$-component zero.

\begin{remark}[Structural mirror of Patch 0551]\label{rem:mirror_of_0551}
Patch~0551 establishes the first-shell perpendicularity $\hat{e}_{ij}^{(S_1)} \cdot \nhat = 0$ for the 30 first-shell-to-first-shell 600-cell edges (the icosahedron edges formed by $S_1(\vhost)$). The present second-shell perpendicularity is the structural mirror: 30 edges instead of 30, but on the second-shell dodecahedron instead of the first-shell icosahedron. Both shells lie in 3-affine-planes orthogonal to $\nhat$ at signed distances $\phi/2$ (first shell) and $1/2$ (second shell) from the origin, so edges within either shell are perpendicular to $\nhat$ by the same chord-difference argument. The structural recurrence reflects the general fact that any shell $S_n(\vhost)$ has all its vertices at the same $\nhat$-projection (by $G_n$ inner-product primitives at progressively deeper levels), so all within-shell edges are perpendicular to $\nhat$. This shell-by-shell perpendicularity pattern is a candidate for future general-shell registration if higher-shell perpendicularity primitives are needed.
\end{remark}

\begin{remark}[On the absence of an icosahedral-transitivity step]\label{rem:no_transitivity_step}
The perpendicularity result Theorem~\ref{thm:g24_hardened} is established for \emph{every} pair $(w_{j_1}, w_{j_2}) \in E_{S_2}(\vhost)$ without needing the $\Icoh$-action explicitly; the chord-difference argument is uniform in the choice of endpoints. This is structurally different from Patch~0588 (host-to-second-shell projection), which used dodecahedral transitivity as a consistency cross-check. The perpendicularity result is independent of pair identity because the $\nhat$-component of any second-shell vertex is the SAME ($1/2$, by G2), not just related by $\Icoh$-action; the chord difference annihilates this common $\nhat$-component identically. The 30-fold uniformity of Theorem~\ref{thm:g24_hardened} is thus a consequence of the SINGLE-VALUE structure of G2 (all 20 second-shell vertices share the same $\nhat$-projection), not of dodecahedral transitivity per se.
\end{remark}

\section{Preliminaries and imported inputs}\label{sec:preliminaries}

We re-use the framework setup of Patch~0587 \S2 without re-derivation: 600-cell edge graph $\Gsixcell$ with 120 vertices on $S^3 \subset \Reals^4$ at unit norm (H2); vertex-aligned Reading~C with $\nhat = \vhost$ (H1); $\Hthree = \Icoh$ residual symmetry at the host vertex (H3) inherited from Patch~0571 Theorem~1.1 + Capotauro~v2.0 FI-C-RC-2 + Patch~0419 Finding~C-W37.

The single imported input is the G2 inner-product primitive:

\begin{lemma}[G2 inner-product primitive; imported from Patch~0587 Theorem~1.1]\label{lem:G2_imported}
Under hypotheses (H1)--(H3), every second-shell vertex $w_j \in S_2(\vhost)$ satisfies
\begin{equation}\label{eq:G2_inner_product}
w_j \cdot \vhost \;=\; \frac{1}{2}
\end{equation}
uniformly across $j \in \{1, \ldots, 20\}$.
\end{lemma}

\begin{proof}[Proof of Lemma~\ref{lem:G2_imported}]
This is Patch~0587 Theorem~1.1, at publication-grade Layer~3 unconditional rigor. Imported here without re-derivation. $\square$
\end{proof}

\subsection{Proof of the edge count Lemma~\ref{lem:30_edges_count}}\label{sec:edge_count_proof}

\begin{proof}[Proof of Lemma~\ref{lem:30_edges_count}]
We establish the edge count by direct enumeration on a canonical representative $w_0 = \tfrac{1}{2}(1,1,1,1) \in S_2(\vhost)$ for $\vhost = (1, 0, 0, 0)$, then propagate via dodecahedral transitivity.

\textbf{Direct enumeration of 600-cell neighbours of $w_0$ in $S_2(\vhost)$.}
A 600-cell vertex $v$ is adjacent to $w_0$ in $\Gsixcell$ iff $v \cdot w_0 = \phi/2$ (the 600-cell first-shell inner-product primitive G1, Patch~0571 Theorem~1.1, applied with host $w_0$). For $v$ in the second coordinate family (the 16 vertices $\tfrac{1}{2}(\pm 1, \pm 1, \pm 1, \pm 1)$): the inner product is $v \cdot w_0 = (1/4) \sum_i \epsilon_i \in \{-1, -1/2, 0, 1/2, 1\}$, none of which equals $\phi/2 \approx 0.809$. Hence no second-family vertex is 600-cell-adjacent to $w_0$.

For $v$ in the third coordinate family (the 96 vertices of even permutations of $(\pm\phi/2, \pm 1/2, \pm 1/(2\phi), 0)$ with independent signs): we seek $v \cdot w_0 = \phi/2$. Each $v$ in the third family has $v \cdot w_0 = (1/2)(\pm a \pm b \pm c)$ where $\{|a|, |b|, |c|\} = \{\phi/2, 1/2, 1/(2\phi)\}$ as a multiset (i.e., $\{a, b, c\}$ is a signed-and-permuted version of $\{\phi/2, 1/2, 1/(2\phi)\}$) and the zero coordinate is absent from the inner product. The inner product range is bounded by $(1/2)(\phi/2 + 1/2 + 1/(2\phi)) = (\phi + 1 + 1/\phi)/4 = (\phi^2 + 1)/(4 \cdot \phi^{-1} \cdot \phi) = \ldots = \phi/2$ at the all-positive choice (using $\phi^{-1} = \phi - 1$ and $\phi^2 = \phi + 1$, the sum $\phi + 1 + 1/\phi = \phi^2 + (\phi - 1)\phi/\phi = 2\phi$, so $(1/2)(2\phi)/2 = \phi/2$).

Thus exactly the all-positive sign choice on the non-zero entries $\{\phi/2, 1/2, 1/(2\phi)\}$ yields $v \cdot w_0 = \phi/2$ for each even permutation that places one of these entries in each of three non-zero positions. The 600-cell-neighbours of $w_0$ from the third family are therefore the 12 vertices obtained from the 12 even permutations of (1,2,3,4) (= $|A_4| = 12$) by applying each permutation to the all-positive base $(\phi/2, 1/2, 1/(2\phi), 0)$.

\textbf{Distribution of these 12 neighbours across shells of $\vhost$.}
For each such neighbour $v$, the first coordinate of $v$ is one of $\phi/2$, $1/2$, $1/(2\phi)$, or $0$ depending on which entry of $(\phi/2, 1/2, 1/(2\phi), 0)$ the even permutation places in position 1. By the shell classification of Patch~0587 \S2.1: first coord $= \phi/2$ means $v \in S_1(\vhost)$; first coord $= 1/2$ means $v \in S_2(\vhost)$; first coord $= 1/(2\phi)$ means $v \in S_3(\vhost)$; first coord $= 0$ means $v \in S_4(\vhost)$ (the equator).

Counting even permutations that place each base-entry in position 1: there are 12 even permutations of (1,2,3,4) and 4 positions, so on average each base-entry occupies position 1 in $12/4 = 3$ even permutations. We verify by explicit enumeration: for each fixed position-1 occupant $b \in \{\phi/2, 1/2, 1/(2\phi), 0\}$, the remaining three entries permute in the remaining three positions according to an EVEN permutation; the number of even permutations of three elements is $3!/2 = 3$. So exactly 3 even permutations place each base-entry in position 1, giving $3 \times 4 = 12$ even permutations as expected.

Therefore $w_0$ has exactly 3 600-cell-neighbours in each of $S_1(\vhost)$, $S_2(\vhost)$, $S_3(\vhost)$, $S_4(\vhost)$. The total $S_1 + S_2 + S_3 + S_4$ count is $3 + 3 + 3 + 3 = 12$, matching the 12-regularity of $\Gsixcell$ exactly with zero second-family contributions.

\textbf{Total edge count $|E_{S_2}(\vhost)| = 30$.}
By dodecahedral transitivity of $\Icoh$ on $S_2(\vhost)$ (Patch~0587 Lemma~2.2), the 3-neighbour-in-$S_2$ count of $w_0$ extends to every $w_j \in S_2(\vhost)$. Total directed $S_2$-$S_2$ adjacency count = $20 \times 3 = 60$. Total undirected edge count = $60 / 2 = 30$.

\textbf{Coincidence with the dodecahedron edges.}
The 30 edges $E_{S_2}(\vhost)$ are exactly the 30 edges of the regular dodecahedron formed by $S_2(\vhost)$ under $\Icoh$: each $w_j$ has 3 $S_2$-adjacent 600-cell-neighbours, matching the 3-regularity of the dodecahedron graph (since the dodecahedron has 20 vertices and 30 edges, vertex-degree $2 \cdot 30 / 20 = 3$). Hence $E_{S_2}(\vhost)$ coincides with the dodecahedron edge set up to identification of the 600-cell-edge graph structure with the dodecahedron graph structure on $S_2(\vhost)$.
$\square$
\end{proof}

\section{Proof of the perpendicularity identity}\label{sec:proof}

\begin{proof}[Proof of Theorem~\ref{thm:g24_hardened}]
Fix any second-shell-to-second-shell 600-cell edge $(w_{j_1}, w_{j_2}) \in E_{S_2}(\vhost)$, with edge-direction unit vector $\hat{e}_{j_1 j_2}$ defined by~\eqref{eq:edge_unit_def}.

By Lemma~\ref{lem:G2_imported} (G2 imported from Patch~0587) applied independently to $w_{j_1}$ and $w_{j_2}$,
\begin{equation*}
w_{j_1} \cdot \vhost \;=\; w_{j_2} \cdot \vhost \;=\; \frac{1}{2}.
\end{equation*}
Therefore
\begin{equation}\label{eq:chord_difference_inner_product}
(w_{j_2} - w_{j_1}) \cdot \vhost \;=\; (w_{j_2} \cdot \vhost) - (w_{j_1} \cdot \vhost) \;=\; \frac{1}{2} - \frac{1}{2} \;=\; 0.
\end{equation}
Under hypothesis (H1), $\nhat = \vhost$, so~\eqref{eq:chord_difference_inner_product} gives $(w_{j_2} - w_{j_1}) \cdot \nhat = 0$.

The edge-direction vector is a positive scalar multiple of $w_{j_2} - w_{j_1}$ (the chord difference normalised by its length); since the inner product with $\nhat$ is linear, this scaling preserves the zero-projection property:
\begin{equation*}
\hat{e}_{j_1 j_2} \cdot \nhat \;=\; \frac{(w_{j_2} - w_{j_1}) \cdot \nhat}{|w_{j_2} - w_{j_1}|} \;=\; \frac{0}{|w_{j_2} - w_{j_1}|} \;=\; 0,
\end{equation*}
where the chord-difference denominator $|w_{j_2} - w_{j_1}| > 0$ (the two endpoints are distinct since they are joined by a non-trivial $\Gsixcell$-edge). This establishes~\eqref{eq:g24_main} for the chosen edge.

The argument applies identically to every edge $(w_{j_1}, w_{j_2}) \in E_{S_2}(\vhost)$ since the G2 import (Lemma~\ref{lem:G2_imported}) gives the same $1/2$ inner-product value at every second-shell vertex. The 30-fold uniformity of the perpendicularity statement is therefore not via $\Icoh$-transitivity per se, but via the single-value structure of G2 (Remark~\ref{rem:no_transitivity_step}). $\square$
\end{proof}

\section{Consequences and corollaries}\label{sec:consequences}

\subsection{Vanishing first-order substrate-rate asymmetry on second-shell-to-second-shell edges}\label{sec:vanishing_rate_asymmetry}

\begin{corollary}[Vanishing first-order Mechanism A rate asymmetry on second-shell-to-second-shell edges]\label{cor:vanishing_rate_asymmetry}
Under (H1)--(H3) and the Mechanism A rate function $r(\hat{e}) = r_0(1 + \delta\,\hat{e} \cdot \nhat)$ (F.1 v1.0~\S4 MA.1), the rate-asymmetry factor on every second-shell-to-second-shell 600-cell edge vanishes at first order in $\delta$:
\begin{equation}\label{eq:vanishing_rate_asymmetry}
r(\hat{e}_{j_1 j_2}) - r(-\hat{e}_{j_1 j_2}) \;=\; 2 r_0 \delta (\hat{e}_{j_1 j_2} \cdot \nhat) \;=\; 0
\end{equation}
for every $(w_{j_1}, w_{j_2}) \in E_{S_2}(\vhost)$ with edge-direction unit vector $\hat{e}_{j_1 j_2}$.
\end{corollary}

\begin{proof}
Immediate from Theorem~\ref{thm:g24_hardened}: $\hat{e}_{j_1 j_2} \cdot \nhat = 0$ kills the first-order $\delta$-coefficient in $r$, yielding $r(\hat{e}_{j_1 j_2}) = r(-\hat{e}_{j_1 j_2}) = r_0$ identically (independent of $\delta$ at first order). $\square$
\end{proof}

\subsection{Inputs for downstream Route C sequence 2 path-class enumeration A6}\label{sec:use_in_a6}

Corollary~\ref{cor:vanishing_rate_asymmetry}'s vanishing of the first-order rate asymmetry on second-shell-to-second-shell edges is a load-bearing simplification for the path-class weight enumeration at $\mathcal{O}(\delta^2)$ in artifact A6 (\texttt{o\_delta\_squared\_path\_class\_weights.tex}). Path classes that traverse a second-shell-to-second-shell edge contribute a factor proportional to the rate asymmetry on that edge; by Corollary~\ref{cor:vanishing_rate_asymmetry}, these contributions vanish at first order in $\delta$ for the second-shell-edge step. The path classes that survive at $\mathcal{O}(\delta^2)$ therefore have their perturbative insertion(s) on non-second-shell-to-second-shell edges. This is exactly parallel to Patch~0551 Corollary~4.1's observation that the first-shell-to-first-shell edges contribute zero rate-asymmetry at first order, restricting the $\mathcal{O}(\delta^1)$ current at $\vhost$ to host-to-first-shell edges in F.1 v1.0~\S7.

\subsection{Cross-sector consistency with Capotauro v2.0 K3-base protection}\label{sec:capotauro_parallel}

The second-shell perpendicularity result has a structural parallel in Capotauro~v2.0~\S5.6 K3-base protection: under vertex-aligned Reading~C with $\nhat = \vhost$, the K3-base substructure of the Capotauro chirality theorem inherits a perpendicular-to-$\nhat$ protection at first order in the spatial-sector substrate-locality perturbation amplitude $\epsilon$. The shared underlying mechanism is the chord-difference argument: any substructure built from vertices at uniform $\nhat$-projection inherits perpendicular-to-$\nhat$ edge directions, giving first-order substrate-rate protection. The K3-base, the W-bracelet (Capotauro~v2.0~\S5--\S6), the first-shell icosahedron (Patch~0551), and the second-shell dodecahedron (the present Patch) are all instances of this general shell-protection pattern.

\section{Scope, framework-locality, and explicit exclusions}\label{sec:scope_and_exclusions}

Theorem~\ref{thm:g24_hardened} is scope-bounded by hypotheses (H1)--(H3) and the imports from Patch~0587. This section enumerates four classes of generalisations that lie outside the present theorem's scope.

\begin{enumerate}[label=\textbf{(E\arabic*)}]
\item \textbf{Higher-shell perpendicularity primitives are NOT addressed.} The within-shell perpendicularity pattern (every within-shell edge is perpendicular to $\nhat$) extends formally to all shells $S_n(\vhost)$ by the chord-difference argument, but each higher-shell extension requires importing the corresponding $G_n$ inner-product primitive at publication-grade Layer~3 unconditional rigor. The present Patch establishes the $n = 2$ case using G2 (Patch~0587). Future $n = 3, 4, \ldots$ extensions would parallel.

\item \textbf{First-shell-to-second-shell edge-direction projections are NOT addressed here.} The 60 first-shell-to-second-shell 600-cell edges have endpoints at \emph{different} $\nhat$-projections ($\phi/2$ vs $1/2$), so the chord-difference argument does NOT give zero; rather, these edges have non-zero $\nhat$-projections that decompose under the $D_{3d}$ first-shell-vertex stabilizer into five orbit classes. This is the content of artifact A4 (\texttt{first\_shell\_second\_shell\_edge\_orbits.tex}; G2.3), a Route~C sequence~2 artifact conditional on both Patch~0587 (G2) and Patch~0588 (G2.1), deferred per single-artifact-per-Patch discipline.

\item \textbf{Non-vertex-aligned Reading C variants are NOT covered.} The theorem requires $\nhat = \vhost$ (hypothesis (H1)). Alternative Reading~C placements would yield different second-shell within-shell edge directions. Registered as OPEN-FP-F1-5.

\item \textbf{Layer 4 derivation of the 600-cell substrate from CPP A1--A11.} The 600-cell polytope and its shell decomposition are taken as polytope-theoretic input \cite{coxeter1973}. A Layer~4 axiomatic derivation is OPEN-FP-F1-2.
\end{enumerate}

\section{Cross-references and consequences}\label{sec:cross_references}

\subsection{Completion of the A3 + A5 parallel pair}\label{sec:parallel_pair_completion}

With Patches~0588 (G2.1 host-to-second-shell projection) and~0589 (G2.4 second-shell perpendicularity) both landed at publication-grade Layer~3 unconditional rigor, the parallel A3 + A5 pair of Route~C sequence~2 unlocked by Patch~0587's G2 hardening is complete. The two parallel artifacts together with G2 (Patch~0587) form a publication-grade Layer~3 unconditional second-shell building-block trio:
\begin{itemize}[leftmargin=2em]
\item Patch~0587 (G2): $w_j \cdot \vhost = 1/2$ uniformly (second-shell inner-product primitive).
\item Patch~0588 (G2.1): $\hat{w}_j \cdot \nhat = -1/2$ uniformly (host-to-second-shell projection).
\item Patch~0589 (G2.4; this paper): $\hat{e}_{j_1 j_2} \cdot \nhat = 0$ uniformly (second-shell-to-second-shell perpendicularity).
\end{itemize}
This second-shell trio is structurally parallel to the first-shell trio: G1 (Patch~0571), G1.1 host-to-first-shell projection (Patch~0552), G1.2 first-shell perpendicularity (Patch~0551). The first-shell trio is the input for the F.1 v1.0 first-order coefficient computation; the second-shell trio (now in hand) is one of the input groups for the upcoming A6 + A7 sequence-2 second-order coefficient computation.

\subsection{Position relative to OPEN-FP-F1-1 Route C sequence 2}\label{sec:open_fp_f1_1_position}

OPEN-FP-F1-1 Route~C sequence~2 (scoping document \texttt{sketches/F1\_o\_delta\_squared\_extension\_scoping.md}~\S5.1) comprises seven artifacts A1--A7. Status at the close of the present Patch:
\begin{itemize}[leftmargin=2em]
\item A1 (structural theorem): CLOSED at Patch~0585.
\item A2 (G2 hardening): CLOSED at Patch~0587.
\item A3 (G2.1 host-to-second-shell projection): CLOSED at Patch~0588 (parallel partner to the present Patch).
\item A5 (G2.4 second-shell perpendicularity; \textbf{this Patch~0589}): CLOSED.
\item A4 (G2.3 first-shell-to-second-shell edge orbit classification): OPEN; conditional on Patches~0587 + 0588, both now in hand. A4 may proceed as the next Route~C sequence~2 closure Patch.
\item A6 (path-class weight enumeration): OPEN; conditional on A4 + the framework-extension hypothesis (H5) at sketch-document Layer~3 per DG-1 sub-option (A) (per Patch~0585 framework setup).
\item A7 (umbrella assembly producing THEO-DSL-5 candidate): OPEN; conditional on A1--A6.
\end{itemize}
The sequence-2 closure budget (4--8 sessions per the original scoping estimate) has now consumed three of those budgetary sessions across Patches~0587 + 0588 + 0589. Remaining budget: 1--5 sessions for A4 + A6 + A7.

\section{Conclusion --- the F.1 Layer 3 hardening sequence advances to eight artifacts}\label{sec:conclusion}

Theorem~\ref{thm:g24_hardened} with the G2 import from Patch~0587 hardens the second-shell-to-second-shell edge-direction perpendicularity identity G2.4 at publication-grade Layer~3 unconditional rigor: $\hat{e}_{j_1 j_2} \cdot \nhat = 0$ identically across all 30 second-shell-to-second-shell 600-cell edges. The proof is essentially a chord-difference calculation using the single-value structure of G2. The four exclusion classes (E1)--(E4) make explicit what is NOT covered.

\paragraph{Position in the F.1 hardened-theorem sequence (now eight artifacts).}
With Patch~0589, the F.1 Dynamical Substrate Law Layer~3 hardened-theorem sequence has eight artifacts:
\begin{itemize}[leftmargin=2em]
\item \textbf{Patches~0550 + 0551 + 0552 + 0571} hardened the $\mathcal{O}(\delta^1)$ first-shell content (the original four-artifact trio + G1 primitive sequence~1 completion).
\item \textbf{Patch~0585} hardened the $\mathcal{O}(\delta^k)$ all-orders parallel-to-$\nhat$ structural theorem (Route~C sequence~1).
\item \textbf{Patch~0587} hardened the G2 second-shell inner-product primitive (Route~C sequence~2 A2).
\item \textbf{Patch~0588} hardened the G2.1 host-to-second-shell projection (Route~C sequence~2 A3).
\item \textbf{Patch~0589 (this paper)} hardens the G2.4 second-shell perpendicularity (Route~C sequence~2 A5).
\end{itemize}

The next sequence-2 Patch is A4 (first-shell-to-second-shell edge orbit classification G2.3, conditional on G2 + G2.1, both now in hand at publication-grade Layer~3 unconditional rigor). After A4, A6 + A7 close the sequence and complete the OPEN-FP-F1-1 coefficient sub-question via the THEO-DSL-5 candidate $\alpha_2$.

\paragraph{Anti-erasure note.}
The present hardening preserves: (a)~the explicit edge-count derivation in Lemma~\ref{lem:30_edges_count} \S\ref{sec:edge_count_proof} (the 30 dodecahedron edges of $S_2(\vhost)$ coincide with the 30 second-shell-to-second-shell 600-cell edges) rather than the implicit appeal to "the dodecahedron has 30 edges" that would silence the polytope-theoretic identification; (b)~Remark~\ref{rem:no_transitivity_step}'s registration that the 30-fold uniformity is via the single-value structure of G2 (not via $\Icoh$-transitivity per se), distinguishing this artifact's proof structure from the Patch~0588 partner's reliance on dodecahedral transitivity; (c)~Remark~\ref{rem:mirror_of_0551}'s registration of the general shell-by-shell perpendicularity pattern as a candidate observation for future higher-shell extension.

\bibliographystyle{plain}
\begin{thebibliography}{99}

\bibitem{coxeter1973}
H.~S.~M.~Coxeter,
\textit{Regular Polytopes},
3rd ed., Dover Publications, New York, 1973.

\bibitem{humphreys1990}
J.~E.~Humphreys,
\textit{Reflection Groups and Coxeter Groups},
Cambridge Studies in Advanced Mathematics 29, Cambridge University Press, 1990.

\bibitem{abshier_first_shell_perpendicularity}
T.~L.~Abshier,
``First-Shell-to-First-Shell Edge-Direction Perpendicularity in the 600-Cell at Vertex-Aligned Reading~C,''
CPP F.1 hardened-theorem artifact, Patch~0551, 24~May~2026, \texttt{hardened\_theorems/first\_shell\_perpendicularity.tex}.

\bibitem{abshier_g2_hardening}
T.~L.~Abshier and Claude Opus 4.7,
``Second-Shell Inner-Product Primitive (G2) in the 600-Cell at Vertex-Aligned Reading~C,''
CPP F.1 hardened-theorem artifact, Patch~0587, 26~May~2026, \texttt{hardened\_theorems/second\_shell\_inner\_product\_primitive.tex}.

\bibitem{abshier_host_second_shell_projection}
T.~L.~Abshier and Claude Opus 4.7,
``Host-to-Second-Shell Uniform Projection (G2.1) in the 600-Cell at Vertex-Aligned Reading~C,''
CPP F.1 hardened-theorem artifact, Patch~0588, 26~May~2026, \texttt{hardened\_theorems/host\_second\_shell\_uniform\_projection.tex}.

\bibitem{abshier_f1_dynamical_substrate_law}
T.~L.~Abshier,
``The Dynamical Substrate Law: Substrate-Locality of DI-Bit Currents at Vertex-Aligned Reading~C in the 600-Cell,''
CPP Flagship Paper Series (F-line), Version 1.0, 24~May~2026, OSF DOI \texttt{10.17605/OSF.IO/JXE8D}.

\end{thebibliography}

% BEGIN GENERATED IDENTIFIER APPENDIX -- do not edit by hand
\clearpage
\section*{Appendix: Programme identifiers used in this paper}
\addcontentsline{toc}{section}{Appendix: Programme identifiers}

\noindent This programme tracks its results, assumptions and unresolved questions under short identifiers, so that a claim and its dependencies can be cited precisely. Prefixes denote the kind of item: \texttt{THEO-} a proved theorem, \texttt{FI-} a foundational input the derivation assumes, \texttt{PRED-} a registered prediction, \texttt{CONJ-} a conjecture, \texttt{OPEN-} an unresolved question, \texttt{PH-} a problem history. Those appearing in this paper are glossed below; no external document is needed to read them.

\begin{description}
  \item[\texttt{FI-C-RC-2}] \hfill \\ The vertex-aligned reading identifying that 4D direction with the host vertex, fixing the local structure as the icosahedral first shell.
  \item[\texttt{OPEN-FP-F1-1}] \hfill \\ Extend the F.1 substrate-locality umbrella theorem (Theorem 7.1) to \$\textbackslash\{\}mathcal\{O\}(\textbackslash\{\}delta\textasciicircum{}2)\$ by deriving the second-shell inner-product and edge-projection identities for the 600-cell, and determine whether the parallel-to-\$\textbackslash\{\}hat\{n\}\$ structure of \$\textbackslash\{\}vec\{J\}\_\{\textbackslash\{\}text\{DI-net\}\}(v\_\{\textbackslash\{\}text\{host\}\})\$ survives at second order. **Resolution**: both sub-questions closed — structural sub-question at publication-grade Layer 3 unconditional (\$\textbackslash\{\}vec\{j\}\_k(v\_\{\textbackslash\{\}text\{host\}\}) \textbackslash\{\}parallel \textbackslash\{\}hat\{n\}\$ at every \$k \textbackslash\{\}geq 1\$ via THEO-DSL-4) + coefficient sub-question at sketch-document Layer 3 under (H5) (\$\textbackslash\{\}alpha\_2 = -9/\textbackslash\{\}phi\textasciicircum{}2 = -9(2-\textbackslash\{\}phi) \textbackslash\{\}approx -3.438\$ via THEO-DSL-5 candidate; ratio \$\textbackslash\{\}alpha\_2/\textbackslash\{\}alpha\_1 = -3/2\$ exactly).
  \item[\texttt{OPEN-FP-F1-2}] \hfill \\ Derive Mechanism A (F.1 framework axioms MA.1 + MA.2: propagation-rate asymmetry \$r(\textbackslash\{\}hat\{e\}) = r\_0(1 + \textbackslash\{\}delta\textbackslash\{\},\textbackslash\{\}hat\{e\}\textbackslash\{\}cdot\textbackslash\{\}hat\{n\})\$ and framework-local current construction at \$\textbackslash\{\}mathcal\{O\}(\textbackslash\{\}delta\textasciicircum{}1)\$) from CPP primitive axioms A1–A11 alone.
  \item[\texttt{OPEN-FP-F1-5}] \hfill \\ Extend or reformulate F.1 Theorems 5.1 (host-first-shell projection), 5.2 (first-shell perpendicularity), and 7.1 (substrate-locality umbrella) under edge-aligned Reading C (\$\textbackslash\{\}hat\{n\}\$ along a 600-cell short-edge direction; residual \$D\_5\$ symmetry at edge midpoint per Patch 0596 Remark 5.1 anti-erasure correction — see below) and face-aligned Reading C (\$\textbackslash\{\}hat\{n\}\$ perpendicular to a triangular face; residual \$C\_s\$ symmetry under vertex-host convention per Patch 0597 Remark 7.1 anti-erasure correction — see below).
  \item[\texttt{THEO-DSL-4}] \hfill \\ \$\textbackslash\{\}mathcal\{O\}(\textbackslash\{\}delta\textasciicircum{}k)\$ Substrate-Current Parallel-to-\$\textbackslash\{\}hat\{n\}\$ Structural Theorem at All Orders \$k \textbackslash\{\}geq 1\$ (**fourth F-line flagship theorem under THEO-DSL-N sub-prefix convention**; publication-grade Layer 3 **unconditional** for the structural sub-question; **first Route C sequence 1 artifact** under the OPEN-FP-F1-1 trajectory; **closes the structural sub-question of OPEN-FP-F1-1** at all orders \$k \textbackslash\{\}geq 1\$ via symmetry-only argument; coefficient \$\textbackslash\{\}alpha\_k\$ at \$k \textbackslash\{\}geq 2\$ remains OPEN as target of Route C sequence 2)
  \item[\texttt{THEO-DSL-5}] \hfill \\ A Dynamical Substrate Law theorem giving a real coefficient in the substrate rate function at the host vertex.
  \item[\texttt{THEO-DSL-N}] \hfill \\ (family) A proposed family registration for the shell-perpendicularity pattern that recurs at every perturbative order.
\end{description}

\subsection*{How we review}

\noindent Results in this programme are checked by an \emph{AI review panel}: several large language models from different vendors are given the same paper and the same fixed protocol, and asked to find errors and raise objections independently. Objections are recorded and either answered in the text or registered as open problems under the identifiers glossed above.

\noindent This is a \emph{verification method the authors apply to their own work}. It is not journal peer review, it is not equivalent to it, and agreement among panel members is not evidence that a result is correct --- only that no member found a specific fault. Where individual models are named in the text, they are named for reproducibility of the method; model versions change, and a reader should treat the named models as a record of what was run rather than as an endorsement.

% END GENERATED IDENTIFIER APPENDIX

\end{document}
